% !TeX root = ../../Book1.tex
% 纲要标题：### 6.2 紧曲面三角剖分的证明
% 纲要位置：几何学计划纲要.md，第 325 行。
% 正文状态：初稿；尚未正式审核。

% 纲要内容开始
% 
% * 平面 Jordan–Schoenflies 定理与简单弧直化
% 
%   在局部平面图中先建立 Jordan–Schoenflies 的平面 Jordan 曲线版本，以及圆盘内适当嵌入简单弧的直化与延拓引理，再处理有限曲线网。三角剖分的主范围为紧、Hausdorff、第二可数的拓扑二维流形，边界允许存在；平面拓扑引理为本章明确的证明任务，不预设曲面已有光滑图册，也不调用高维三角剖分结论。
% 
% * 交叠边界的局部有限交叉
% * 由曲线网得到有限胞腔分解
% * 圆盘胞腔的三角化
% * 同痕后的公共细分与 Euler 示性数不变性
% 
% 纲要内容结束

\section{紧曲面三角剖分的证明}
\label{g1:sec:0602}

形成有限三角剖分，需要把坐标图中的曲线整理成有限的边与顶点。简单闭曲线可能处处没有切线，圆盘边界也可能无限相交；平面中的分离与直化结论使我们能够逐个调整边界，同时保留圆盘覆盖。平面部分采用有限图与逐次加密的构造；这条证明思路见\cite[第2--4节]{Thomassen1992}。

\subsection{平面 Jordan–Schoenflies 定理与简单弧直化}
\label{g1:subsec:0602:01}

这里的简单弧是 \([0,1]\) 的嵌入像，\term[Jordanquxian]{Jordan 曲线}[Jordan curve]是 \(S^1\) 的嵌入像。\footnote{Camille Jordan（1838--1922），法国数学家，研究群论与拓扑。}嵌入的参数映射可以非常不规则；整个论证只用连续性和单射性。先给出后面所需的平面分离与有限图结论。

\begin{lemma}[有限折线与平面图]\label{g1:lem:polygonal-plane-tools}
	下列事实成立。
	\begin{enumerate}
		\item 平面连通开集中的任意两点可以用一条简单折线连接。
		\item 简单闭折线把平面分成两个连通开集，其中一个有界；二者的边界都是这条折线。闭折线内部的一条简单折线弧若恰在两个端点与边界相交，就把内部切成两个区域，其边界分别由该弧与原边界的两段组成。
		\item 任意有限图在平面中的嵌入都可以重新画成边为折线的嵌入。完全二部图 \(K_{3,3}\) 没有平面嵌入；把一个平面图中的互不相交的有限树各收缩成一点，仍得到平面图。
		\item 若一个有限平面图连通、没有环边、没有割点，且至少有三个顶点（若原图有环边，先将环边细分，再检验这些条件），则它可以从任一圈开始，逐次加入两端在旧图上、内部与旧图不交的简单弧得到。若各边为折线，它的每个面以一个简单圈为边界。
	\end{enumerate}
\end{lemma}

\begin{proof}
	对于第一项，固定起点，把能用有限折线到达的点组成集合。小圆盘内可以添一条线段，所以这个集合及其补集在给定开集中都开；连通性说明它等于整个开集。有限条线段的交点和重叠端点只有有限个，把它们作为顶点得到有限图，删去回路便得到简单折线。

	对于第二项，先给出折线两侧的具体含义。在每条边的内部取窄矩形，在每个顶点取足够小的圆盘，使不相邻的这些小邻域不相交；顶点圆盘内只有两条相接的半径方向，可以沿这两条边的两侧分别绕过顶点。沿整个多边形走一周，得到两条互不相交的窄通路。平面补集中的任一点可以沿折线走向边界，取第一次接触以前的一小段，就到达这两条通路之一；若接触顶点，改在顶点圆盘中的相应扇形内行走。因此补集至多有两个分支。

	再证明两侧不能相通。从不在边界上的点向无穷远发一条射线，先让射线避开顶点且不平行于边，记录它与边界的交点数模 \(2\)。改变射线方向经过一个顶点时，两次交点同时出现或消失，或者一次交点换到相邻边，故奇偶性不变。沿不碰边界的短线段移动起点也不改变奇偶性：在扫过的窄四边形中，每段边界从边界进入还必须离开，交点变化为偶数。有限折线连接把这一结论推广到同一分支中的任意两点。跨过一条边的内部时奇偶性恰好改变，因此确有两个分支。足够大的圆外连通，故只有一个分支无界；窄通路同时说明边界上的每一点都在两分支的闭包中。对于内部弧，把它与边界的两段分别组成简单闭折线，重复奇偶性论证；弧的两侧不同，而每一侧都沿对应边界窄通路连通，恰好得到所说的两个区域。

	第三项的重新画法不要求原边可微。先给环边添一个顶点，使每条边有不同的端点。在每个顶点周围取互不相交的小圆盘，只截到与该顶点相接的边。对每条边取最后离开起点圆盘以后、随后第一次进入终点圆盘以前的闭子弧；这些子弧两两不交，因而有两两不交的开邻域。用第一项在各邻域内改画简单折线，并在顶点圆盘内用互不相交的径向线段接到顶点。\(K_{3,3}\) 包含一个六边圈和连接三对相对顶点的三条弦。在折线画法中，至少两条弦落在该圈的同一侧；它们的四个端点沿圈交错，与第二项的切分结论矛盾。有限折线树的细邻域可以由顶点小圆盘和边的细矩形组成；从叶边开始依次添加矩形和端点圆盘，每次仍是一个闭盘。在该闭盘内把树及其附着边改画成一颗星，就实现树的收缩。这也说明若一个图经过收缩互不相交的树能得到 \(K_{3,3}\)，原图便不能平面嵌入。

	最后证明第四项。已有一个没有割点的真子图 \(H\)。若剩余某条边的两端都在 \(H\) 中，可直接加入。否则，取 \(G\setminus H\) 的一个连通分支；它必须接到 \(H\) 的至少两个不同顶点，因为只有一个接点会成为 \(G\) 的割点。在该分支中选连接两个接点的最短路径，便得到所需的新弧。加入它仍没有割点，有限次便遍历全部边。从初始圈出发，每加入一条弧，它都位于某个旧面的内部；第二项说明它把该面分成两个以简单圈为边界的面。于是归纳得到最后的结论。
\end{proof}

下面两步把折线结论推广到任意连续弧和闭曲线。关键区别是：任意简单弧仍然不能隔断平面，而闭曲线确实会隔断。

\begin{lemma}[简单弧不分离平面]\label{g1:lem:arc-nonseparation}
	若 \(A\subset\mathbb R^2\) 是简单弧，则 \(\mathbb R^2\setminus A\) 道路连通。
\end{lemma}

\begin{proof}
	固定 \(x,y\notin A\)，取 \(d>0\)，使它们到 \(A\) 的距离都大于 \(4d\)。按参数把 \(A\) 分成有限段 \(A_1,\ldots,A_k\)，每段直径小于 \(d\)。不相邻的闭段互不相交，所以其间距离的最小值为正；取 \(0<\delta<d\) 小于这些距离。若只有一两段，直接任取这样的 \(\delta\)。在每段上取更细的点列，使相邻点间的子弧直径小于 \(\delta/8\)，并在每个点处画边长 \(\delta/2\)、边平行于坐标轴的正方形。记第 \(i\) 段所用正方形边界的并为 \(G_i\)，把有限交点添作顶点。

	相邻正方形的边界至少有两个公共点；由\cref{g1:lem:polygonal-plane-tools}最后一项，或直接删去任一顶点检查，\(G_i\) 没有割点。相邻的 \(G_i,G_{i+1}\) 包含同一个端点正方形；不相邻的图不交。每一对相邻图都包含在以公共弧端点为中心、半径 \(3d\) 的圆盘内，所以 \(x,y\) 位于这对图的无界面。

	需要一个有限图的观察：这种链状排列的图若把一点围在有界面内，必有一对相邻图已经把它围住。证明如下。由前一引理，整个并没有割点，各有界面的边界都是圈。若观察不成立，取围住该点的圈 \(C\)，使它所用图的最大下标 \(j\) 最小，再使 \(C\) 使用 \(G_j\setminus G_{j-1}\) 的边数最少。若 \(C\) 只在 \(G_{j-1}\cup G_j\) 内，已经矛盾；否则 \(C\) 在 \(G_{j-1}\) 上至少有两段。\(G_{j-1}\) 没有割点，故可取一条内部避开 \(C\) 的路径连接其中两段：在 \(G_{j-1}\) 中取连接两段的最短路径，并截去首次接触之前和末次接触之后的部分即可。它把 \(C\) 切成两个圈；点不在新路径上，必在其中一个圈的内部。该圈或者降低最大下标，或者减少所数的边，均与选择矛盾。

	因此 \(x,y\) 都位于 \(G=\bigcup_iG_i\) 的无界面。这个面是连通开集，可用简单折线连接两点。每一段 \(A_i\) 都被相应正方形覆盖，而每个正方形内部与 \(G\) 的无界面不交；故所得折线避开 \(A\)。
\end{proof}

\begin{theorem}[Jordan 曲线分离定理]\label{g1:thm:jordan-separation}
	任意 Jordan 曲线 \(C\subset\mathbb R^2\) 的补集恰有两个分支：有界分支 \(D\) 和无界分支 \(E\)，且 \(\partial D=\partial E=C\)。两分支都是道路连通的开集。
\end{theorem}

\begin{proof}
	先证至少有两个分支。取两条竖直支撑线，使 \(C\) 位于它们之间，在每条支撑线上取 \(C\) 的最高点 \(p_1,p_2\)。曲线分成从 \(p_1\) 到 \(p_2\) 的两段。在两支撑线之间再取一条竖线，这两段与该线都相交；沿该线找到两个相邻的、分别属于两段的交点 \(a,b\)，使开线段 \((a,b)\) 避开 \(C\)。这是可行的，因为两组交点是互不相交的紧集。另从 \(p_1\) 向上、经曲线上方的水平线、再降到 \(p_2\)，画一条内部避开 \(C\) 的折线 \(L\)。若线段 \((a,b)\) 的中点可以在补集中到达无穷远，它就可以接到 \(L\) 的内部；截取首次接触，所得连接弧的内部同时避开 \(C\)、\([a,b]\) 和 \(L\)。于是 \(C\)、\([a,b]\)、\(L\) 及这条连接弧组成一个 \(K_{3,3}\) 的细分：两类顶点分别是 \(p_1,p_2\) 与连接弧在 \([a,b]\) 上的端点，以及 \(a,b\) 与连接弧在 \(L\) 上的端点。这与\cref{g1:lem:polygonal-plane-tools}矛盾。因此补集不连通。

	对补集的任意分支 \(U\)，考察从 \(U\) 内出发、恰在终点接触 \(C\) 的折线。这些折线所能到达的端点在 \(C\) 中稠密。事实上，任取一个开子弧 \(J\subset C\)，\(C\setminus J\) 是简单闭弧；由\cref{g1:lem:arc-nonseparation}，可用避开该闭弧的折线连接 \(U\) 中的一点和另一个分支中的一点。它第一次离开 \(U\) 时只能碰到 \(J\)，这就给出一个所需的可达点。子弧可以任意缩小，稠密性随之成立。

	假设补集有三个分支 \(U_1,U_2,U_3\)，在曲线上选三段互不相交的短闭子弧 \(J_1,J_2,J_3\)。每个 \(U_i\) 中可以画一棵有限的折线树，分别接到三个 \(J_j\)，且除这些端点外避开 \(C\)：先用上一段选可达点，再在连通开集 \(U_i\) 内连接三条到达折线，删除多余回路。三棵树互不相交。把各树及每个 \(J_j\) 上连接三个接触点的最短子弧各收缩为一点，得到 \(K_{3,3}\)，仍与有限图引理矛盾。所以恰有两个分支。

	足够大的圆外是连通的，因此恰有一个无界分支。补集是开集，平面的连通开集道路连通。边界只能落在 \(C\) 上，而上述可达点在 \(C\) 上的稠密性说明 \(C\) 上每一点都属于每个分支的边界。
\end{proof}

\begin{lemma}[横切弧与有限弦网]\label{g1:lem:jordan-crosscuts}
	设 \(C\) 是 Jordan 曲线，\(D\) 是其有界补分支。若简单折线弧 \(P\) 的内部在 \(D\) 中，且两个端点是 \(C\) 上的不同点，那么 \(P\) 把 \(D\) 切成两个区域，边界分别由 \(P\) 与 \(C\) 的两段组成。设 \(G\) 为由 \(C\) 和 \(D\) 中有限条折线弧组成的连通有限图，将环边先作细分，并假设细分后的图没有割点。则 \(G\) 的每个有界面都以 Jordan 曲线为边界；在某个面中添加一条上述弧，恰好把该面分成两个。
\end{lemma}

\begin{proof}
	记 \(C\) 的两段为 \(A_1,A_2\)，\(C_i=P\cup A_i\) 的有界内部为 \(D_i\)。原曲线的无界外部 \(E\) 是避开 \(C_i\) 的连通无界集，故在 \(C_i\) 的外部。\(A_2\) 的内点属于 \(\overline E\) 又不在 \(C_1\) 上，因此也在 \(C_1\) 的外部；同理 \(A_1\) 的内点在 \(C_2\) 外部。于是 \(D_1\) 不碰 \(C\)，且不在 \(E\) 中，故 \(D_1\subset D\)，同理 \(D_2\subset D\)。

	\(D_1\) 连通且不碰 \(C_2\)，所以整体位于 \(C_2\) 的一侧。若 \(D_1\subset D_2\)，闭包包含关系将使 \(A_1\) 的内点落在 \(\overline{D_2}\)，与上一段矛盾。因此 \(D_1\cap D_2=\varnothing\)。在 \(P\) 的任一内点处取足够小的圆盘，使 \(P\) 在其中是一条折线直径；若恰为折点，先用扇形坐标把两段边变直。在这个小盘中，\(C_1,C_2\) 都只有同一段 \(P\)。Jordan 定理说明每个 \(D_i\) 都接近该点；两侧半盘连通，因此其中一侧属于 \(D_i\)。两内部互不相交，说明它们恰好占据相反的两侧。

	所以 \(D_1\cup P^\circ\cup D_2\) 在 \(D\) 中是开集。它也是 \(D\) 中的闭集，因为它等于 \((\overline{D_1}\cup\overline{D_2})\cap D\)。\(D\) 连通且这个集合非空，故它等于 \(D\)。这就证明 \(D\setminus P=D_1\sqcup D_2\)，并且各自边界正是 \(P\cup A_i\)。有限图部分再按\cref{g1:lem:polygonal-plane-tools}的逐次加弧构造归纳；这里每次只调用刚证明的横切弧结论。
\end{proof}

把 \(D\cup C\) 实现为闭圆盘，需要将边界上的对应延拓到内部。Jordan--Schoenflies 定理的构造同时细化原像和像中的网格，使两侧相邻点的控制在极限中都保留下来，从而得到连续且互逆的映射；\cref{g1:fig:double-mesh}给出构造的示意。\footnote{Arthur Sch\"onflies（1853--1928），德国数学家，研究拓扑与晶体群。}

% !TeX root = ../Book1.tex
\begin{figure}[htbp]
	\centering
	\begin{tikzpicture}[scale=0.88, every node/.style={font=\small}]
		\begin{scope}
			\draw[thick] (-2,-1.4) .. controls (-2.5,-0.2) and (-2.1,1.6) .. (-0.8,1.8) .. controls (0.2,1.4) and (1.1,2.1) .. (1.9,1.3) .. controls (2.4,0.3) and (1.8,-1.6) .. (0.9,-1.8) .. controls (-0.3,-1.4) and (-1.2,-2) .. cycle;
			\draw[GeometryBlue,thick] (-1.2,-0.9) rectangle (1.2,0.9);
			\foreach \x in {-0.6,0,0.6}{\draw[GeometryBlue] (\x,-0.9)--(\x,0.9);}
			\foreach \y in {-0.45,0,0.45}{\draw[GeometryBlue] (-1.2,\y)--(1.2,\y);}
			\draw (-1.2,0.9)--(-0.8,1.8);
			\draw (1.2,-0.9)--(0.9,-1.8);
			\node at (0,2.25) {原区域 \(D\)};
			\node[fill=white,inner sep=2pt] at (0,-0.2) {\(K_n\)};
			\node at (-1.8,1.7) {\(C\)};
			\node at (0,-2.3) {紧集附近的网格逐次变细};
		\end{scope}
		\draw[->,thick] (2.7,0.2)--(3.7,0.2);
		\draw[<-,thick] (2.7,-0.2)--(3.7,-0.2);
		\begin{scope}[xshift=6.4cm]
			\draw[thick] (-1.8,-1.8) rectangle (1.8,1.8);
			\draw[GeometryBlue,thick] (-1.2,-0.9) rectangle (1.2,0.9);
			\foreach \x in {-0.6,0,0.6}{\draw[GeometryBlue] (\x,-0.9)--(\x,0.9);}
			\foreach \y in {-0.45,0,0.45}{\draw[GeometryBlue] (-1.2,\y)--(1.2,\y);}
			\draw (-1.2,0.9)--(-1.8,1.8);
			\draw (1.2,-0.9)--(1.8,-1.8);
			\node at (0,2.25) {目标正方形 \(Q\)};
			\node at (0,-2.3) {再加密像侧，并把新弦复制回去};
		\end{scope}
	\end{tikzpicture}
	\caption{双向网格控制的示意。每次只在内部添加有限条弦；原边界上的对应保持不变。两边都要加密，才能保证极限及其逆都连续。图中的光滑边界只为示意，证明允许任意 Jordan 曲线。}
	\label{g1:fig:double-mesh}
\end{figure}


\begin{theorem}[Jordan--Schoenflies 定理]\label{g1:thm:jordan-schoenflies}
	任意两条平面 Jordan 曲线之间的同胚都可以延拓为整个平面的同胚。特别地，每条 Jordan 曲线及其有界补分支的并同胚于闭圆盘，而且给定的边界同胚可以一并延拓。
\end{theorem}

\begin{proof}
	先把给定边界同胚写成 \(f:C\to\partial Q\)，其中 \(Q\) 是正方形；若这个情形成立，对两条曲线分别构造到正方形的延拓再复合即可。记 \(D\) 为 \(C\) 的内部。我们构造一列有限弦网 \(G_n\subset\overline D\)、\(G'_n\subset Q\)，始终包含边界，内部边为折线。每个图没有割点，其删去边界后的部分连通；两图具有相同的顶点、边、面关系，边界对应始终是 \(f\)。在这个阶段只固定顶点的对应，不固定内部边上每一点的参数。

	从 \(D\) 可沿折线到达的边界点中选可数稠密集 \(A\)，这是\cref{g1:thm:jordan-separation}证明中得到的。取两个这样的点，用 \(D\) 中的简单折线连接，得到初始弦；在正方形内取对应弦。于是得到 \(G_0,G'_0\)。以后依次把 \(A\) 的各点加入顶点集：从新点沿到达折线走进 \(D\)，再连到旧图内部，截到第一次碰到旧图为止。它在某个旧面内成为一条新弦。在对应的正方形面内也可画出相应弦；该面内部连通，端点位于折线边或已有可达的边界点，先取端点处的小段，再用连通开集中的简单折线相接即可。\cref{g1:lem:jordan-crosscuts}保证两面以相同方式各分为两面。

	同时必须加密网格。令
	\[
		K_n=\{x\in\overline D:\operatorname{dist}(x,C)\ge 1/n\}.
	\]
	在 \(D\) 内加入有限个小正方形的边及其中的横、竖网格，使 \(K_n\) 被这些正方形的内部覆盖，且网格的每格直径小于 \(1/n\)。所有这些边与旧图的内部折线只有有限个交点，并与 \(C\) 不交。必要时添加有限条内部折线，使新图没有割点且内部连通。这个补充可以如下完成：先连接各分支；若有割点 \(v\)，在 \(v\) 两侧的分支中各取靠近 \(v\) 的点，在 \(D\setminus\{v\}\) 中用折线连接。连通开集删去一点仍连通，因为一条经过该点的折线可以在一个小圆盘内绕过它。把新交点作为顶点后，这条连接把割点处分开的块合并，使块数至少减少一；按有限图的块数归纳即可结束。每次新增弧都可拆成有限次在旧面内加弦，因而可以在 \(Q\) 中逐条复制其组合关系。

	接着反过来，在 \(Q\) 的内部加入足够细的有限网格，使到边界距离至少 \(1/n\) 的部分被直径小于 \(1/n\) 的网格包围；再把这些弦逐条复制回 \(D\)。新端点若落在已有边的内部，就在对应边上选一个新点，保持已有各点的次序；若在边界上，则选取 \(f(A)\) 中的点。可在边界附近略弯折网格末端以做到这一点，内部网格不受影响。仍用前一段的加弦过程补足连通性。后加的弦只细分已有面，不会破坏先前的网格。因此，对于原像与像两侧各自的每个内部紧集，以及每个 \(\varepsilon>0\)，从某一步起，紧集上每一点都位于一个由图中边组成的简单小圈内部，而这个小圈连同内部的直径小于 \(\varepsilon\)。若点恰在网格边或顶点上，就取环绕它的两行、三行小格的外框；把网格取得更细即可达到同一要求。

	把所有顶点及整个边界组成集合 \(F\)，相应集合记为 \(F'\)。构造给出双射 \(f_0:F\to F'\)，且两集合分别在 \(\overline D,Q\) 中稠密。它保持有限弦网中每个圈的内外关系：这在每次加弦时由横切弧引理保持。现在证明上述对应延拓为同胚。

	取 \(x\in D\)。在某个图中选围住 \(x\) 的内部圈 \(L\)，其对应圈 \(L'\) 也在 \(Q\) 内部。趋近 \(x\) 的 \(F\) 中点最终在 \(L\) 内，像点因而留在 \(L'\) 围出的紧区域中。任取其像点的两个聚点 \(y,z\)。若 \(y\ne z\)，利用像侧越来越细的网格，在 \(y\) 附近取两个互不相交的嵌套网格圈，它们都把 \(y\) 与 \(z\) 分开。对应的两个原像圈也互不相交，至多一个经过 \(x\)。取另一个圈 \(M\)；趋近 \(x\) 的点最终都在 \(M\) 的同一侧，像点也都在对应圈的同一侧，不可能同时趋近位于相反侧且不在圈上的 \(y,z\)。故像点只有一个极限，定义它为 \(h(x)\)。同样的论证对反向顶点对应成立，得到 \(k:Q^\circ\to D\)。这里聚点不会落在边界：环绕任一内部点的内部小圈及其对应圈已经把聚点限制在内部的紧区域。常值顶点序列也是允许的趋近序列，所以延拓在每个已有顶点上与 \(f_0\) 一致。

	这些延拓连续。给定 \(h(x)\) 的一个小邻域，取像侧两个互不相交的嵌套内部小圈，都把 \(h(x)\) 围在其中，并令闭内部落在该邻域；至少一个对应原像圈不经过 \(x\)，选取这一个。它必须围住 \(x\)，否则稠密顶点的对应违反内外关系。这个原像圈的内部便是 \(x\) 的一个邻域，整个邻域的像都落在给定小邻域中。同理 \(k\) 连续。两个复合映射在稠密顶点集上都是恒等映射，连续性给出 \(k\circ h=\mathrm{id}\)、\(h\circ k=\mathrm{id}\)。

	最后检查边界。设 \(x_m\in D\) 趋近 \(a\in C\)，而某个像子列趋近 \(b\in Q\)。若 \(b\in Q^\circ\)，内部圈及其原像紧区域迫使 \(x_m\) 留在 \(D\) 的一个紧子集，矛盾。故 \(b\in\partial Q\)。若 \(b\ne f(a)\)，在 \(C\) 上选两个 \(A\) 中的点 \(a_1,a_2\)，把 \(a\) 与 \(f^{-1}(b)\) 分开。它们最终都已加入图；因图的内部连通，可以在图内找到一条从 \(a_1\) 到 \(a_2\)、其余部分均在内部的简单弦。该弦及其对应弦把两边各分成两块。靠近 \(a\) 的点位于与 \(a\) 相邻的块，其像位于与 \(f(a)\) 相邻的块，后者的闭包不含 \(b\)，矛盾。因此 \(b=f(a)\)。紧性说明整个像序列都趋近 \(f(a)\)，从而 \(h\cup f\) 连续。它是紧空间到 Hausdorff 空间的连续双射，故为闭盘同胚。

	还要延拓到外部。先选各曲线内部的一点作为原点，对两平面作反演 \(x\mapsto x/\|x\|^2\)。每条曲线的外部连同无穷远点变成一条 Jordan 曲线的闭内部，原来的无穷远点对应原点。已经证明的闭盘延拓给出反演后两内部之间的同胚。可使它把原点映到原点：先把目标闭盘映到正方形，在正方形中从给定内部点与目标内部点分别向四角作扇形三角剖分，逐三角形的仿射映射给出固定边界、移动该点的同胚，再共轭回去。于是反演回去得到外部同胚，与内部同胚在原边界上相合。两部分的粘合连续，其逆也同样连续，故得到全平面同胚。
\end{proof}

\begin{lemma}[简单弧的直化与延拓]\label{g1:lem:proper-arc-straightening}
	平面简单弧的每个内点都有一个开矩形型邻域，使该邻域与原弧的交集在某个坐标下恰为水平直径。若 \(B\) 是闭 Jordan 盘，\(A\subset B\) 是简单弧，且 \(A\cap\partial B\) 恰为其两个端点，则 \(B\setminus A\) 有两个分支，闭包都是闭 Jordan 盘。给定把这两个端点送到另一闭盘中一条直径端点的边界同胚，可以延拓为把 \(A\) 送到该直径的闭盘同胚。
\end{lemma}

\begin{proof}
	先证局部结论。在所给内点 \(p\) 两侧各选一个短子弧。每个短子弧内都有从原弧补集沿线段可达的点：取靠近该子弧、但不在原弧上的点，再在原弧上选取离它最近的点，连接这两个点的开线段不碰原弧；起点取得足够近，最近点仍在指定短子弧内。补集由\cref{g1:lem:arc-nonseparation}连通，所以可用其中的折线连接两条到达线段，删去回路，得到一条与原弧只在所选两端点相交的简单弧。它与经过 \(p\) 的原子弧组成 Jordan 曲线。\cref{g1:thm:jordan-schoenflies}把这条曲线变成圆；再在圆上的对应点附近取普通矩形坐标，并缩小矩形以避开原弧剩下的紧集，就得到所求直化。

	对闭盘中的弧，将它分别与盘边界的两段组成 Jordan 曲线。重复\cref{g1:lem:jordan-crosscuts}的两个有界内部的论证：两内部不交且包含在 \(B\) 中；本引理刚证的局部直化使 \(A\) 每个内点附近都有两个半矩形，两个内部各占一侧；它们与 \(A^\circ\) 的并在 \(B^\circ\) 中同时开闭。因此它们恰为 \(B\setminus A\) 的两块，闭包只沿 \(A\) 相交。任取从 \(A\) 到目标直径、与给定两个端点相容的同胚。它与给定的边界同胚在两个子盘的边界上分别组成同胚；对两子盘应用\cref{g1:thm:jordan-schoenflies}并沿 \(A\) 粘合，便得到整个闭盘的同胚。
\end{proof}

这里有一个以后会反复用到的直接推论：位于一个平面图中的闭 Jordan 盘，其边界有双侧环带邻域。用全平面同胚把边界变成单位圆，再取 \(1-\varepsilon<r<1+\varepsilon\) 的逆像即可。有限条在一点相交、除此之外互不相交的弧也可以同时直化为有限条径向线段：先把两条组成的弧直化，再在各半盘内依次用上一引理直化其余弧；每次缩小中心邻域以避开各弧已经截去的紧尾部。因弧数有限，此过程结束后得到一个顶点小圆盘和互不相交的边矩形。因此，有限嵌入图具有由顶点盘、边带组成的细邻域。

\subsection{交叠边界的局部有限交叉}
\label{g1:subsec:0602:02}

现在回到曲面。紧性允许选有限个坐标圆盘，但不能保证它们的边界相交有限次；我们要改变边界，仍保留原来较小圆盘的覆盖性质。下面的论证只在一个坐标图内进行，因此同时适用于曲面上的有限条曲线。

\begin{lemma}[边界的有限化]\label{g1:lem:disk-boundary-regularization}
	设曲面 \(M\) 中有有限条简单曲线组成的嵌入有限图 \(G\)，以及一个坐标图内的三个闭 Jordan 盘
	\[
		B_0\subset\operatorname{int}B_1\subset B_1\subset\operatorname{int}B_2.
	\]
	存在闭 Jordan 盘 \(B\)，使 \(B_0\subset\operatorname{int}B\subset B\subset\operatorname{int}B_2\)，且 \(\partial B\) 与 \(G\) 只有有限个交点。可以进一步使交点均不为 \(G\) 的顶点，并且在局部矩形中两条边分别为横、竖直径。
\end{lemma}

\begin{proof}
	先把 \(B_2\) 当作平面中的圆盘。在 \(G\cap\operatorname{int}B_2\) 中，称两端在 \(\partial B_2\)、内部在盘内的一段边为穿入段。虽然穿入段可能无限多，碰到 \(B_1\) 的穿入段只有有限多：每条这样的段都从 \(\partial B_2\) 走到与它正距离的 \(B_1\)，而一条边的参数映射一致连续，所以它占据的参数区间长度有统一正下界；这些参数区间互不相交，因此只能有有限多个。先在 \(G\) 的有限个顶点处用小盘截断，也把这些顶点附近的有限段一并列入。

	把这些有限段与 \(\partial B_2\) 合在一起，必要时在盘内添有限条弧，使所得有限图没有割点。\cref{g1:lem:proper-arc-straightening}及\cref{g1:thm:jordan-schoenflies}允许固定 \(\partial B_2\)，把这个有限图改画成折线图：按逐次加弧的次序，第一次延拓边界同胚，以后在当前弧所在的两个子盘内分别延拓，再粘合。每一步保持已经画好的边不动。

	在这个坐标变化以后，未列入的穿入段仍不碰新的 \(B_1\)。沿 \(\partial B_1\) 取有限个小正方形，闭正方形都避开 \(B_0\) 及这些未列入段，并使它们覆盖 \(\partial B_1\)。为何能同时避开全部未列入段？在 \(\partial B_1\) 的每一点，若它们有点列趋近该点，则原有限边中的一个子列有参数极限；该极限落在 \(B_1\) 上，所处穿入段也应列入，且一个小参数邻域属于这同一穿入段，矛盾。必要时先把 \(B_1\) 略缩到仍包含 \(B_0\) 的内层盘，便可使闭方形与所有未列入段分离。

	这些正方形连同它们围住的 \(B_0\) 构成有限折线区域；取围住 \(B_0\) 的外边界，并在有限个自接触顶点处用很小的折线绕开，得到 \(B_0\) 与 \(\partial B_2\) 之间的一条简单闭折线。它只可能碰到已列入的有限折线段，所以交点有限。最后在每个交点附近作小的折线改动，避开顶点、重叠和相切：沿每段取不平行于所遇旧边的小线段，并把端点避开有限条直线即可。改动仍在所选小方形内，故既不破坏包含关系，也不会碰到未列入的段。把新边界映回原坐标便得结论。
\end{proof}

局部横交的意思只是两条弧在一个小矩形中呈十字交叉，并没有涉及导数。对于紧无边界曲面，先取有限个坐标盘 \(B_{0,i}\) 的内部覆盖曲面，并为每个盘预留两层较大的坐标盘。逐个应用\cref{g1:lem:disk-boundary-regularization}，已选边界组成旧图，便得到仍覆盖曲面的有限个闭盘 \(B_i\)，其边界两两有限横交，且没有三条边界交于一点。避开有限个已有交点就能满足最后一项。

\subsection{由曲线网得到有限胞腔分解}
\label{g1:subsec:0602:03}

上述圆盘边界形成有限嵌入图 \(G\)。在每条没有交点的圆周上任添一个顶点，就可以把它也看成有限条边。\(M\setminus G\) 的每个分支都位于某个坐标盘的内部：选该分支中的一点和包含它的盘；在不越过该盘边界的情况下，连通分支不可能从盘内走到盘外。这里使用的是盘边界在本地坐标中已经证明的分离性质。

一个分支的边界不一定已经是一个简单圈：同一顶点或同一边可能在边界行走中出现多次，分支也可能有若干个洞。沿图取细邻域，可以把重复的边界行走拆开，使余下区域具有互不相交的边界圆周。

\begin{lemma}[有限曲线网的圆盘化]\label{g1:lem:finite-network-disk-cells}
	设紧无边界曲面 \(M\) 上的有限嵌入图 \(G\) 的每个补分支都包含在某个平面坐标图内。则可以添加有限条简单弧，并对所得曲线网作有限细分，使 \(M\) 被分成有限个闭 Jordan 盘；不同盘只沿整条边或顶点相交。
\end{lemma}

\begin{proof}
	按\cref{g1:lem:proper-arc-straightening}后的构造，为每个顶点取互不相交的小盘，为每条截短的边取细矩形，各矩形只在两个端部接到相应顶点盘。边数和顶点数有限，紧的不相交部分可以用互不相交的邻域隔开，因此各次缩小能够同时完成。所得闭邻域 \(N\) 的边界由有限段弧拼成有限条互不相交的圆周。顶点盘与边矩形本身已经是有限个盘形块。

	\(M\setminus\operatorname{int}N\) 的每个分支 \(R\) 是一个紧平面区域，其边界是上述圆周中的有限条。这样的分支只有有限多个：每个分支都接到 \(\partial N\) 的至少一条圆周，而一条圆周的外侧局部属于唯一分支。若有分支没有边界，它在 \(M\) 中同时开闭；在所考虑的连通分支内这会迫使 \(G\) 为空，可先添一个坐标小圆周排除该情形。每个 \(R\) 的边界恰由一个外圈和有限个内圈组成。由 Jordan 分离定理，任一边界圈的一侧属于 \(R\)，另一侧不属于；在最外圈内反复排除内圈的闭内部，就得到 \(R\)，不存在额外遗漏区域，否则其边界还应包含另一条 \(\partial R\) 圈。

	若 \(R\) 有内圈，利用其内部的折线连通性，从一个内圈接一条简单弧到外圈，只在两端碰边界。取该弧的窄矩形带并切开，相当于把这两个边界圈合并成一个圈，内圈数减少一。弧与边界相接处使用半矩形坐标，故切开后仍是带边界的紧曲面；在平面中也可直接把矩形带从区域移去，它的两条长边取代原弧。有限次以后只剩一个边界圈，由\cref{g1:thm:jordan-schoenflies}所得区域是闭盘。把先前取出的矩形带作为单独的闭盘形块加回来，便得到有限盘形分解。最后将各端点和交点统一添为顶点，使两盘共有的边都成为完整边。
\end{proof}

\subsection{圆盘胞腔的三角化}
\label{g1:subsec:0602:04}

盘形分解已将曲面分成有限个闭圆盘，各块沿完整的边或顶点相交。在每个圆盘中添中心并连至边界分点，可以把它分成三角形；公共边上使用同一组分点，使这些三角形能够拼成全曲面的剖分。

\begin{definition}[有限三角剖分]\label{g1:def:surface-triangulation}
	曲面 \(M\) 的一个\term[sanjiaopoufen]{有限三角剖分}[finite triangulation]是有限单纯复形 \(K\) 及同胚 \(|K|\to M\)。这里 \(K\) 由有限个顶点、边和三角形组成，每个单形的面仍在 \(K\) 中，两个单形的交集是各自的一个公共面或空集；\(|K|\) 表示按这些面粘合后的拓扑空间。若 \(M\) 有边界，还要求 \(\partial M\) 是一些边和顶点的并。
\end{definition}

\begin{theorem}[紧拓扑曲面的有限三角剖分]\label{g1:thm:surface-triangulation}
	每个紧、Hausdorff、第二可数的二维拓扑流形都有有限三角剖分；允许存在边界。对于给定的、与边界相容的有限嵌入图，可以预先指定有限个点，并使该图和这些点成为剖分的子复形；这里与边界相容指图与边界的交集是它的若干整条边及顶点。尤其可以保留有限条两两不相交的内嵌简单闭曲线和适当嵌入简单弧。
\end{theorem}

\begin{proof}
	先考虑无边界情形。有限圆盘覆盖及其边界正则化给出\cref{g1:lem:finite-network-disk-cells}所需曲线网，因而得到有限闭盘分解。若预先指定有限条曲线与弧，就从这些曲线组成的图开始添加圆盘边界；在交点或预先指定点处分边，后续步骤保持它们作为图的一部分。给每个盘的边界添足至少三个顶点，选一个内部点。把该盘同胚到正多边形，在内部点与边界各顶点之间连线，再映回来，就把该盘分成有限个三角形。盘与盘之间的边界划分事先统一，所以这些三角形可以沿相同的边粘合。

	这一步可能只得到一种容许重边的三角形粘合，例如两条不同的边具有相同端点。作一次重心细分便可消除此问题：给每个原顶点、原边和原盘分别赋予不同的新顶点，每条严格包含链“顶点属于边，边属于盘”对应一个小三角形。为避免一个盘边界中同一边重复导致重复链，使用上一引理的闭 Jordan 盘分解，其边界已经单射，且共有部分均为完整的边或顶点。于是一个包含链决定唯一小三角形，交集恰为共同的子链，真正得到有限单纯复形。同胚在每个闭盘上连续并相容，有限闭集粘合给出全局连续双射；紧性和 Hausdorff 性使它成为同胚。

	若 \(M\) 有边界，先把两个副本沿边界逐点粘合，得到双倍曲面 \(\widehat M\)。\cref{g1:prop:boundary-invariance}保证这里的边界是内在定义的；边界点的半平面图把两个半圆盘粘成整圆盘，因此 \(\widehat M\) 是无边界曲面。它紧且第二可数；Hausdorff 性也可直接在图中验证：不同内点用各自副本的邻域分离，两个不同边界点先在 \(M\) 中取闭包不交的小半圆盘，再把对应的两半邻域粘合。由\cref{g1:prop:boundary-circles}，\(\partial M\) 是有限个圆周的不交并，在双倍中成为有限条简单闭曲线。

	对 \(\widehat M\) 应用无边界情形，并把这些边界圆周以及预先给定弧的两个副本加入初始曲线网，指定点也保留。所得剖分以 \(\partial M\) 为子复形；每个小三角形内部不越过这条子复形，所以选取原副本一侧的三角形及其面，就得到 \(M\) 的有限三角剖分。若 \(M\) 不连通，局部连通性使各分支开，紧性说明分支只有有限个，分别构造再取不交并即可。
\end{proof}

紧性在这里提供有限覆盖、有限边界分支与最终的有限剖分。第二可数性是本书对曲面的统一假设；在紧 Hausdorff、局部 Euclid 的情形中，它也可由有限图册推出：每个坐标图有可数基，把有限个图的基合并即可。

\subsection{同痕后的公共细分与 Euler 示性数不变性}
\label{g1:subsec:0602:05}

两个任意拓扑剖分可能有无限交叉：一条边可以在另一条边附近来回摆动无穷多次，而两边仍各自是简单弧。先用环境同痕调整一份剖分，使两套边网只有有限横交，再逐个细分这些交点，就能寻找共同的三角剖分。

\begin{definition}[环境同痕]\label{g1:def:ambient-isotopy}
	空间 \(M\) 上的\term[huanjingtonghen]{环境同痕}[ambient isotopy]是一族同胚 \(H_t:M\to M\)，其中 \(0\le t\le1\)，\(H_0=\mathrm{id}_M\)，并且 \((x,t)\mapsto(H_t(x),t)\) 是 \(M\times[0,1]\) 的同胚。它连续移动整个空间，所以同时控制被移动图形与其周围的部分。对本节的紧 Hausdorff 曲面，只要 \((x,t)\mapsto H_t(x)\) 连续且每个 \(H_t\) 是同胚，上述乘积映射便是紧空间到 Hausdorff 空间的连续双射，逆映射自动连续。
\end{definition}

\begin{lemma}[有限图的小环境同痕]\label{g1:lem:graph-small-isotopy}
	设 \(K\) 是紧曲面 \(M\) 的一个有限三角剖分，\(G\subset M\) 是另一个与边界相容的有限嵌入图，即 \(G\cap\partial M\) 是 \(G\) 的若干整条边与顶点。存在从恒等映射开始的环境同痕 \(H_t:M\to M\)，使 \(H_1(G)\) 的每条边在 \(K\) 的三角形坐标中由有限条线段组成。可以把同痕限制在 \(G\) 的任意给定邻域内；有边界时保持 \(\partial M\) 作为集合不变，并可保持预先共有的边界顶点不动。
\end{lemma}

\begin{proof}
	取 \(G\) 的互不相交顶点小盘及边的窄矩形，使用\cref{g1:lem:proper-arc-straightening}后的有限图邻域构造，并使这些邻域都落在指定邻域中。先在每个顶点盘内把相交的有限条初始边改成径向星形，再把截短边作为各矩形的中线处理。由于 \(K\) 有限，在 \(|K|\) 中充分细分以后，可以在每个盘、矩形内画出一颗折线星或一条简单折线中线，与原来的端点及循环次序一致：沿原弧取有限条小三角形链，连接相邻小邻域中的点，删除回路；不同边使用互不相交的窄矩形，故不会产生新交叉。顶点附近则在一张三角形扇形的有限平面图中按给定循环次序选择短线段，再用小折线连接到边矩形。

	旧星形与新星形之间的对应可以延拓为固定顶点盘边界的同胚。逐条应用闭盘适当弧的延拓引理，把一个扇区分成两个，再在各剩余子盘内延拓即可；只有有限个扇区。各边矩形中，固定矩形边界，将旧中线送到新折线的同胚同样由\cref{g1:lem:proper-arc-straightening}给出。按“先顶点盘、后截短矩形”的次序复合这些同胚，它们在各自支撑边界上恒等，故与外部的恒等映射粘合。

	每个固定闭盘边界的同胚都同痕于恒等映射。把盘识别为单位圆盘，若该同胚为 \(u\)，定义
	\[
		U_t(x)=
		\begin{cases}
			t\,u(x/t),&\|x\|\le t,\quad t>0,\\
			x,&\|x\|\ge t,
		\end{cases}
		\qquad U_0(x)=x.
	\]
	在 \(\|x\|=t\) 两式相等；每个 \(U_t\) 的逆由同一公式将 \(u\) 换成 \(u^{-1}\) 得到，且 \(\|U_t(x)-x\|\le2t\) 保证在 \(t=0\) 连续。这就把有限次盘内修改逐一提升为环境同痕。边界半盘的修改保持其直径上的点或指定端点，反射成整盘后使用相同构造，再限制回半盘即可。
\end{proof}

\begin{theorem}[同痕后的公共细分]\label{g1:thm:common-subdivision-isotopy}
	设 \(K,L\) 是紧拓扑曲面 \(M\) 的两个有限三角剖分。存在环境同痕 \(H_t\)，使 \(K\) 与 \(H_1(L)\) 有一个有限公共细分。有边界时，可以保持边界为集合。
\end{theorem}

\begin{proof}
	对 \(L\) 的边网应用\cref{g1:lem:graph-small-isotopy}，并在有限个小盘内再作小的折线调整，使它避开 \(K\) 的顶点、没有重叠边段，且与 \(K\) 的边只作有限次横交；边界上允许共有边段，把所有端点统一细分即可。两个有限折线网的并是有限嵌入图。每个补分支都在 \(K\) 的一个三角形及 \(H_1(L)\) 的一个三角形之内，但可能有洞或重复的边界行走；按\cref{g1:lem:finite-network-disk-cells}的顶点盘、边带及连接内圈的构造，把它进一步切成闭盘，再在每个盘内三角化。取出的顶点盘和边带也分别按两个原剖分的边切开，因而每个新小三角形都包含在两个剖分各自的某个单形内。统一共有边上的分点，最后作重心细分，即得有限公共细分。
\end{proof}

\begin{proposition}[Euler 示性数与剖分无关]\label{g1:prop:euler-triangulation-independent}
	若紧曲面 \(M\) 的有限三角剖分有 \(V\) 个顶点、\(E\) 条边、\(F\) 个三角形，则整数
	\[
	\chi(M)=V-E+F
	\]
	不依赖所选三角剖分。更一般地，任何有限闭盘分解的顶点数减边数加盘数也等于 \(\chi(M)\)。这个数称为\term[Eulershu]{Euler 示性数}[Euler characteristic]。\footnote{Leonhard Euler（1707--1783），瑞士数学家，在分析与几何等领域作出广泛贡献。}
\end{proposition}

\begin{proof}
	先记住三种局部计数。给一条边加一个顶点时，\(V,E\) 各增加 \(1\)；用一条新弧把一个盘分成两个盘时，\(E,F\) 各增加 \(1\)；在一个盘内部加一个顶点并向其边界的 \(r\) 个顶点连线时，\(V,E,F\) 分别增加 \(1,r,r-1\)。三种变化都保持 \(V-E+F\)。

	任意有限细分也保持这个数。先在各旧边上加入所有新顶点，再处理每个旧盘内部的有限连通边网。把它连同盘边界画在平面中，选择一个生成树；树有 \(V'-1\) 条边，且补集只有一个面，因为可从叶边开始，用小圆盘绕过叶边，并逐条删去叶边。每增加一条非树边，它与旧图中连接两端的一条简单路径组成 Jordan 曲线。这条新边的两侧因该曲线而不能再相通；它内部有矩形邻域，所以只把所在的一个面分成两个，其他面不变。因此有界面数为 \(E'-(V'-1)\)，即这个闭盘分解满足 \(V'-E'+F'=1\)。边界上若有 \(b\) 个顶点，就恰有 \(b\) 条边，它们在交替和中抵消，故盘内新增的顶点、边与面取代原来一个盘时仍贡献 \(1\)。对各旧盘求和，便证明细分前后计数相同，也同时证明盘形分解与其三角化的计数相同。

	给两个三角剖分应用\cref{g1:thm:common-subdivision-isotopy}。环境同痕只移动单形，不改变它们的数目；同痕后的两个剖分具有公共细分，而细分保持计数，故两者的 \(V-E+F\) 相等。
\end{proof}

Euler 示性数的 \(V-E+F\) 计数以盘形面为单位。环面上有一条圆周，其补集是一个环带，这个环带就不是一个可以按“一个面”计数的闭盘。后文使用多边形或带分解计算 Euler 示性数时，会先把各块细分为本节允许的盘形块。
